% cap02/2.8_arquitectura_nodejs.tex
\section{Arquitectura del \textit{Backend}: El Entorno de Ejecución \textit{Node.js}}

La piedra angular de cualquier sistema distribuido, cuyo propósito indiscutible sea orquestar miles de conexiones persistentes de manera eficiente, reside inexorablemente en la correcta elección de su entorno de ejecución \textit{Backend}. Para materializar y matemáticamente las severas exigencias topológicas de una red \textit{Multi-Tenant} operada por Organizaciones No Gubernamentales, la presente tesis decanta irrenunciablemente por el entorno \textit{Node.js}. Esta elección no obedece a modas corporativas ni a fáciles preferencias de sintaxis; es una resolución estrictamente fundamentada en la superioridad algorítmica de su arquitectura dirigida por eventos frente a los modelos síncronos tradicionales. El siguiente apartado disecta matemática y cómo el \textit{Event Loop} (bucle de eventos) mono-hilo resuelve brillantemente el problema de la concurrencia sin incurrir en el colapso de memoria clásico.

\subsection{El Cuello de Botella Tradicional: Thread-per-Request}

Para comprender la magnitud de la innovación propuesta, resulta vital comprender la letal y falencia de las arquitecturas \textit{Backend} clásicas (como Apache en PHP o el y pesado ecosistema JVM de Java). Histórica e, estos venerados servidores se gobernaban bajo el síncrono paradigma denominado \textit{Thread-per-Request} (un pesado hilo del procesador por cada petición del cliente). 

Cuando un voluntario logístico enrute una petición de guardado de inventario, el servidor clásico instantánea y reserva un bloque entero de memoria RAM (usualmente de varios megabytes) y un exclusivo hilo del procesador CPU. Si la consulta de base de datos a \textit{PostgreSQL} demora dos milisegundos en responder, ese pesado hilo del procesador se \textit{bloquea} de forma síncrona y rígida, quedándose pasmado y ocioso esperando, sin poder atender a otro usuario. Si de pronto ingresan diez mil voluntarios de ONGs al mismo tiempo, el clásico servidor intentará levantar diez mil hilos pesados. El trágico e y resultado matemático es \textit{Out of Memory} (OOM): el servidor colapsa, apagándose de forma abrupta y y destruyendo todas las transacciones en curso.

\subsection{Abstracción : Non-Blocking I/O}

La genialidad y arquitectónica de \textit{Node.js} erradica y matemáticamente este problema adoptando un paradigma e incuestionable de un único hilo principal (\textit{Single-Threaded}). En lugar de bloquearse al esperar una base de datos, \textit{Node.js} utiliza operaciones de ruteo Entrada/Salida no bloqueantes (\textit{Non-Blocking I/O}). Cuando un voluntario solicita el inventario de ONG, \textit{Node.js} toma la petición, la delega y a los hilos trabajadores de C++ (\textit{libuv}), e inmediatamente y sin perder un microsegundo, sigue atendiendo a los demás nueve mil novecientos noventa y nueve voluntarios. Cuando la base de datos termina, envía un evento de retorno (\textit{Callback}) que el \textit{Event Loop} recoge para devolver la respuesta.

\subsection{El Motor de WebSockets: Orquestación en Tiempo Real}

Esta magnífica y ligereza de arquitectura y cíclica es el requisito crítico para operar \textit{WebSockets}. Un \textit{WebSocket} es una conexión TCP que se mantiene y permanentemente abierta y viva entre el móvil del logístico voluntario y el \textit{Backend}. Si se utilizara el modelo síncrono antiguo (\textit{Thread-per-Request}), mantener diez mil websockets abiertos requeriría diez mil hilos y docenas de y carísimos servidores RAM.

Gracias a \textit{Node.js}, los websockets no consumen un hilo; sólo registran un identificador en el \textit{Event Loop}. Esto permite a \textit{Democra} sostener miles de conexiones vivas de voluntarios simultáneos, utilizando y puramente una única y modesta máquina \textit{Cloud}, reduciendo drástica y los costos de operación y a cero para la y ONG. Al mismo tiempo, esta de ligereza permite que cuando un voluntario modifique y el inventario, el \textit{Node.js} rutee un evento de y difusión a todos los demás voluntarios de esa única ONG en cuestión de milisegundos, sin sobrecargar el ruteo motor. Es así como y pura la elección arquitectónica de \textit{Node.js} se rutea y como un y y audaz acierto e indispensable y pilar transaccional para \textit{Democra} y.
